\section{函数列与函数项级数 IV}

\subsection{幂级数 continue}

\subsubsection{收敛半径判定}

求幂级数的收敛半径有类似比值判别法和根值判别法的方法。

\begin{theorem}[Cauchy-Hadamard 定理]
    设 $A = \limsup_{n \to \infty} \sqrt[n]{|a_n|}$, 则幂级数 $\sum_{n=0}^\infty a_n (x-x_0)^n$ 的收敛半径
    $$
    R = \begin{dcases}
        0 & ,\ A = \infty \\
        +\infty & ,\ A = 0 \\
        \frac{1}{A} & ,\ 0 < A < \infty
    \end{dcases}
    $$
\end{theorem}

\begin{theorem}[比值判别法]
    设 $A = \lim_{n \to \infty} \abs{\frac{a_{n+1}}{a_n}}$, 则幂级数 $\sum_{n=0}^\infty a_n (x-x_0)^n$ 的收敛半径
    $$
    R = \begin{dcases}
        0 & ,\ A = \infty \\
        +\infty & ,\ A = 0 \\
        \frac{1}{A} & ,\ 0 < A < \infty
    \end{dcases}
    $$
\end{theorem}

\begin{theorem}[Abel 第二定理]
    设 $\sum\limits_{n = 0}^{\infty} a_n x^{n}$ 的收敛半径为 $R\ (0 < R < +\infty)$，那么
    \begin{itemize}
        \item 若 $\sum\limits_{n = 0}^{\infty} a_n R^{n}$ 收敛，则 $\sum\limits_{n = 0}^{\infty} a_n x^{n}$ 在 $[0, R]$ 上一致收敛；
        \item 若 $\sum\limits_{n = 0}^{\infty} a_n (-R)^{n}$ 收敛，则 $\sum\limits_{n = 0}^{\infty} a_n x^{n}$ 在 $[-R, 0]$ 上一致收敛。
    \end{itemize}

    \begin{proof}
        仅证明第一条。对于 $x \in [0, R]$ 有
        $$
        \sum_{n = 0}^{\infty} a_n x^{n} = \sum_{n = 0}^{\infty} \ab(a_n R^{n}) \ab(\frac{x}{R})^{n}
        $$
        因为 $\sum\limits_{n = 0}^{\infty} a_n R^{n}$ 收敛，$\ab(\frac{x}{R})^{n}$ 单调且一致有界，根据 Abel 判别法得证。
    \end{proof}
\end{theorem}

\subsubsection{幂级数连续性、可积性、可导性}

\begin{theorem}[连续性定理]
    设幂级数 $\sum\limits_{n = 0}^{\infty} a_n x^{n}$ 的收敛半径为 $R$。那么
    \begin{itemize}
        \item 在 $(-R, R)$ 上连续；
        \item 在 $(-R, R)$ 上可积，且
        $$ 
        \int_a^b \sum_{n = 0}^{\infty} a_n x^{n} \dd{x} = \sum_{n = 0}^{\infty} \int_a^b a_n x^{n} \dd{x} = \sum_{n = 0}^{\infty} \frac{a_n}{n + 1} \ab(b^{n + 1} - a^{n + 1})
        $$ 
        \item 在 $(-R, R)$ 上可导，且
        $$ 
        \ab(\sum_{n = 0}^{\infty} a_n x^{n})' = \sum_{n = 0}^{\infty} \ab(a_n x^{n})' = \sum_{n = 1}^{\infty} n a_n x^{n - 1}
        $$ 
    \end{itemize}
\end{theorem}

\subsubsection{函数的幂级数展开}

对 Taylor 展开进行扩展，可以得到函数的幂级数展开。

\begin{definition}
    对于函数 $f(x)$ 若在定义域上有
    $$
    f(x) = \sum_{n = 0}^{\infty} a_n (x - x_0)^{n}
    $$
    那么称 $f(x)$ 是\textbf{可展开}的。
\end{definition}

若 $f(x)$ 有任意阶导数，那么根据 Taylor 展开可以得到函数的 Taylor 级数：
$$
f(x) = \sum_{n = 0}^{\infty} \frac{f^{(n)}(x_0)}{n!} (x - x_0)^{n}
$$

但是 Taylor 级数不一定收敛，收敛也不一定和函数本身相等。例如
$$
f(x) = \begin{dcases}
    \E^{-\frac{1}{x^2}} & ,\ x > 0 \\
    0 & ,\ x \le 0
\end{dcases}
$$
